\documentclass{beamer}
\usetheme{metropolis}

\usepackage{amsmath}
\usepackage{listings}
\usepackage{hyperref}
% Define Sauder green
\definecolor{saudergreen}{RGB}{91,191,33}
% Or: \definecolor{saudergreen}{HTML}{A3BD00}
% Apply it to the frame title background and text
\setbeamercolor{frametitle}{bg=saudergreen, fg=white}

\begin{document}

\title{PIBC: Think like a housing economist}
\author{Tom Davidoff\\Sauder School of Business, UBC}
\frame{\titlepage}

\begin{frame}
	\frametitle{The time value of money}
	\begin{itemize}
		\item A dollar now generally viewed as more valuable than a dollar in the future
			\begin{itemize}
				\item Inflation
				\item Can invest the dollar now and earn a return
				\item Impatience 
			\end{itemize}
		\item A discount rate translates future to present dollars
			\begin{itemize}
				\item We will call this rate $d$
				\item The ``present value'' of a payment of \$$x$ in $t$ years is $\frac{x}{\left[1+d\right]^t}$.
				\item Example if discount rate is 10\% \$1,000 in 7 years has PV
					$$\frac{1000}{\left[1.10\right]^7} \approx \$513$$
			\end{itemize}
		\item Project Net Present Value (``NPV'') 
			\begin{itemize}
				\item Is sum of PV associated with project at all dates
				\item Including negative of investment today (date 0)
					$$NPV = \sum_{t=0}^{\infty}x_{t}\left[1+d\right]^{-t}$$
				\item (assuming constant discount rate ``flat yield curve'')
			\end{itemize}
	\end{itemize}
\end{frame}

\begin{frame}
	\frametitle{More on discounting}
	\begin{itemize}
		\item Arguably, parents job to inculcate low discount rate
		\item IRR: ``internal rate of return''
			\begin{itemize}
				\item Like NPV Jeopardy
				\item ``At this discount rate, NPV of project is zero''
				\item Positive IRR only if $IRR>d$
			\end{itemize}
		\item Project pro formas
			\begin{itemize}
				\item All the info needed to compute NPV and IRR
				\item Much more detail than just dates and cash flows
				\item Important ratio: ``cap rate''
					\begin{itemize}
						\item Net operating income (NOI) in year 1 divided by cost\\
								Acquisition cost for simple purchase\\
								Total project cost o/w
					\end{itemize}
				\item NOI vs cashflow
					\begin{itemize}
						\item NOI ignores financing (interest and amortization)
						\item Cashflow $\approx$ NOI after financing costs
					\end{itemize}
			\end{itemize}
	\end{itemize}
\end{frame}

\begin{frame}
	\frametitle{The world's simplest pro forma}
	\begin{itemize}
		\item Buy land for \$1M today
		\item Spend \$1M building something also today
		\item Sell for \$2.25M in one year
		\item Discount rate: 10\%
		\item NPV: -1M - 1M + $\frac{2.5M}{1.1} \approx \$272K$
		\item IRR: -2M + 2.5M/(1+IRR) = 0 $\Rightarrow$ IRR=25\%
	\end{itemize}
\end{frame}

\begin{frame}
	\frametitle{Useful game to understand discounting}
	\begin{itemize}
		\item Pick your favourite musician's last name
			\begin{itemize}
				\item A-M: $d=3\%$
				\item N-Z: $d=10\%$
			\end{itemize}
		\item Match up with a partner
		\item Try to create a loan that has positive NPV to both of you
		\item Loan has amount $L$, interest rate $r$, and term of 1 year
		\item Borrower NPV: $L\left[1-\frac{1+r}{1+d}\right]$
			\begin{itemize}
				\item Positive iff $r<d$
			\end{itemize}
		\item Lender NPV: $L\left[\frac{1+r}{1+d}-1\right]$
			\begin{itemize}
				\item Positive iff $r>d$
			\end{itemize}
		\item If same discount rate, impossible to have +NPV for both
	\end{itemize}
\end{frame}

\begin{frame}
	\frametitle{Segue: from game to why afforable housing mandates are dicey}
	\begin{itemize}
		\item What do markets do, ideally?
		\item Get everyone to agree on ratios of benefits to cost
		\item I have a discount rate of 10\%, you 5\%
			\begin{itemize}
				\item You should lend me money at a rate between 5\% and 10\%
				\item Positive NPV to both of us
			\end{itemize}
		\item Very similar: 
			\begin{itemize}
				\item You can build an apartment for \$500K
				\item I value the apartment at \$550K
			\end{itemize}
		\item With ``complete markets'' trades happen until agreement
			\begin{itemize}
				\item I have so much housing, I no longer value more than your cost
			\end{itemize}
		\item Affordability mandates:
			\begin{itemize}
				\item Someone willing to pay \$400K 
				\item Market price \$600K
				\item Force sale at, say, \$350K
				\item Buyer gets +\$50K, seller gets -\$250K
				\item ``Pareto gain'': seller gives buyer \$150K!
			\end{itemize}
	\end{itemize}
\end{frame}

\begin{frame}
	\frametitle{The money value of time and urban land markets}
	\begin{itemize}
		\item ``Location, location, location''
		\item The classical ``monocentric city model''
			\begin{itemize}
				\item Everyone commutes to location 0: the CBD
				\item Cost of commuting: $t$ per unit of distance
				\item How does land price change with distance?
			\end{itemize}
		\item ``Bid-rent curves''
			\begin{itemize}
				\item Willingness to pay for land at distance $x$ is $R(x)$
				\item To be indifferent between living at $x$ a tiny bit further
				\item Consume $l(x)$ units of land
				\item $R(x)l(x) = R(x+dx)l(x+dx) + t\cdot dx$
				\item So slope of $R$ curve is $-\frac{t}{l(x)}$
				\item Land valuation falls with distance, 
					\begin{itemize}
						\item But more slowly if you consume more land
						\item Hence steeper curve closer to downtown
						\item (also $t$ bigger downtown)
					\end{itemize}
			\end{itemize}
	\end{itemize}
\end{frame}

\begin{frame}
	\frametitle{Equilibrium with multiple uses}
	\begin{itemize}
		\item Remember, slope of bid-rent curve is $-\frac{t}{l(x)}$
		\item Closer to downtown will be steeper curve
	\end{itemize}
	\pgfimage[height=6cm]{monocentric}
\end{frame}

\begin{frame}
	\frametitle{Who should be closer to town given bid rent curve logic?}
	\begin{itemize}
		\item Remember, slope of bid-rent curve is $-\frac{t}{l(x)}$
		\item Starbucks vs WalMart
		\item Yuppie versus family with 1 earner, 2 kids, 2 pets?
		\item Apartment building versus house with yard?
		\item Office building versus warehouse?
	\end{itemize}
\end{frame}

\begin{frame}
	\frametitle{Cap rates, elasticity of supply, and equilibrium\\Where should cap rate be lower?\\Inelastically supplied city if demand is growing}
	\pgfimage[height=6cm]{elasticity}
\end{frame}
	
	
	
		
\end{document}
